\newpage
\section*{Solucionario completo}

\noindent Cada respuesta incluye la clave oficial y su justificación analítica o matemática detallada.

\subsection*{Bloque A. Fundamentos teóricos y conceptuales}

\sol{1}{A}{\textbf{Tema:} Definición de externalidad. \textbf{Explicación:} Una externalidad existe cuando las decisiones de producción o consumo de un agente impactan directamente en la función de utilidad o de costos de terceros sin mediar transacción de mercado ni compensación pecuniaria, desconectando el costo privado del costo social.}

\sol{2}{B}{\textbf{Tema:} Costo marginal social. \textbf{Explicación:} El costo marginal social incluye la totalidad de recursos utilizados por la sociedad: $CMgS = CMgP + CMgE$. Al ser la externalidad negativa, $CMgE > 0$, por lo que $CMgS > CMgP$.}

\sol{3}{C}{\textbf{Tema:} Sobreproducción en externalidad negativa. \textbf{Explicación:} Como las empresas solo consideran sus costos marginales privados ($P = CMgP$), ignoran el daño externo marginal ($CMgE$), produciendo en exceso respecto al nivel donde el beneficio social iguala al costo social ($P = CMgS$).}

\sol{4}{D}{\textbf{Tema:} Externalidad positiva en consumo. \textbf{Explicación:} En externalidades positivas de consumo, los individuos solo toman en cuenta su satisfacción personal ($BMgP$). Al no internalizar el beneficio aportado a los demás ($BMgE > 0$), la demanda privada queda por debajo de la social y el mercado produce menos del óptimo social ($Q_C < Q^*$).}

\sol{5}{A}{\textbf{Tema:} Impuesto pigouviano. \textbf{Explicación:} Arthur Pigou demostró que para internalizar la externalidad, el impuesto unitario debe igualar el daño marginal que sufre la sociedad en el punto de eficiencia: $t^* = CMgE(Q^*)$, alineando la condición privada $P = CMgP + t^*$ con la social $P = CMgS$.}

\sol{6}{B}{\textbf{Tema:} Subsidio pigouviano. \textbf{Explicación:} El subsidio pigouviano reduce el costo percibido por el agente o incrementa su beneficio marginal privado, desplazando el equilibrio hacia la cantidad socialmente óptima donde $BMgS = CMg$.}

\sol{7}{C}{\textbf{Tema:} Supuestos del Teorema de Coase. \textbf{Explicación:} Ronald Coase (1960) postuló que si los derechos de propiedad están asignados sin ambigüedad y no existen fricciones ni costos de negociación (costos de transacción iguales a cero), las partes negociarán hasta alcanzar el óptimo de Pareto.}

\sol{8}{D}{\textbf{Tema:} Independencia distributiva en Coase. \textbf{Explicación:} Bajo ausencia de efectos riqueza, el nivel final de producción y daño ambiental alcanzado tras la negociación es idéntico independientemente de quién tenga el derecho (contaminar o estar libre de contaminación), cambiando únicamente los pagos compensatorios entre las partes.}

\sol{9}{A}{\textbf{Tema:} Límites del Teorema de Coase. \textbf{Explicación:} Cuando los agentes afectados son numerosos y dispersos, los costos de coordinar, negociar y supervisar el cumplimiento son prohibitivos (altos costos de transacción) y surgen conductas estratégicas oportunistas (free-riding y holdout), impidiendo el acuerdo voluntario.}

\sol{10}{B}{\textbf{Tema:} Permisos transferibles (Cap-and-Trade). \textbf{Explicación:} El regulador controla la cantidad agregada admisible emitiendo un número fijo de licencias transferibles. El libre comercio de derechos genera un precio de mercado que induce a reducir contaminación a aquellas firmas con menores costos marginales de abatimiento.}

\sol{11}{C}{\textbf{Tema:} Eficiencia en costos de instrumentos económicos. \textbf{Explicación:} Cuando las empresas tienen costos de reducción distintos, los estándares uniformes obligan a esfuerzos costosos a firmas ineficientes. Los instrumentos de mercado descentralizan la decisión: quien limpia barato abate más y quien limpia caro compra permisos o paga el impuesto, minimizando el costo total social.}

\sol{12}{D}{\textbf{Tema:} Propiedades del bien público puro. \textbf{Explicación:} Un bien público puro es no rival (el consumo de un individuo no reduce la cantidad ni calidad disponible para los demás) y no excluible (es técnica o económicamente inviable impedir que alguien lo consuma una vez provisto).}

\sol{13}{A}{\textbf{Tema:} Clasificación de recursos comunes. \textbf{Explicación:} Los recursos comunes (como los peces en alta mar o el agua de un acuífero subterráneo) son rivales (el agua que extrae un usuario no la puede usar otro) pero no excluibles (cualquiera con acceso físico puede perforar o extraer), lo que provoca sobreexplotación.}

\sol{14}{B}{\textbf{Tema:} Problema del polizón (Free-Rider). \textbf{Explicación:} Dado que nadie puede ser excluido del beneficio del bien público una vez creado, cada individuo racional tiene incentivos a fingir desinterés o subdeclarar su valoración para no pagar su cuota, esperando beneficiarse de la contribución voluntaria de los demás.}

\sol{15}{C}{\textbf{Tema:} Condición de Samuelson. \textbf{Explicación:} Paul Samuelson (1954) demostró que debido a la no rivalidad, una unidad adicional del bien público beneficia a todos los consumidores simultáneamente. Por ende, la suma vertical de los beneficios marginales (disposiciones a pagar) debe igualar el costo marginal social: $\sum RMS_i = RMT$.}

\sol{16}{D}{\textbf{Tema:} Suma vertical de demandas. \textbf{Explicación:} Para bienes privados rivales se suman cantidades horizontalmente a cada precio ($Q = \sum q_i$). Para bienes públicos no rivales, todos consumen la misma cantidad provista ($G$), por lo que se suman verticalmente los precios de reserva: $P(G) = \sum P_i(G)$.}

\sol{17}{A}{\textbf{Tema:} Equilibrio de Lindahl. \textbf{Explicación:} Erik Lindahl propuso que cada individuo pague una fracción fiscal del costo total igual a su valoración marginal individual ($t_i = RMS_i$). Si los agentes revelaran sus preferencias con veracidad, los ingresos cubrirían exactamente el costo del bien público óptimo.}

\sol{18}{B}{\textbf{Tema:} Mecanismo VCG. \textbf{Explicación:} Bajo el mecanismo VCG, declarar la verdadera valoración es una estrategia estrictamente dominante para cada participante: si la declaración de un agente cambia la decisión social, paga una tasa que refleja la pérdida neta impuesta a los demás participantes.}

\sol{19}{C}{\textbf{Tema:} Tragedia de los comunes. \textbf{Explicación:} Bajo libre acceso, la entrada de usuarios prosigue hasta que el ingreso medio iguala el costo marginal del esfuerzo (disipando la renta económica a cero). Como el usuario no considera el efecto de su extracción sobre el rendimiento del resto, la explotación sobrepasa ampliamente el óptimo social.}

\sol{20}{D}{\textbf{Tema:} Bienes meritorios de Musgrave. \textbf{Explicación:} En bienes meritorios (como el agua potable clorada, la educación primaria y las vacunas), la sociedad considera que toda persona debe acceder a ellos independientemente de su capacidad o deseo inicial de pago, corrigiendo fallas de información o preferencias miopes.}

\sol{21}{A}{\textbf{Tema:} Selección adversa vs Riesgo moral. \textbf{Explicación:} La selección adversa ocurre antes de contratar (un agente posee características privadas no observables, como su estado de salud o calidad del producto). El riesgo moral ocurre después de contratar (el agente toma decisiones no observables de esfuerzo o cuidado que afectan el riesgo).}

\sol{22}{B}{\textbf{Tema:} Modelo de Lemons de Akerlof. \textbf{Explicación:} Los compradores solo están dispuestos a pagar por la calidad esperada promedio. A ese precio, los dueños de autos buenos ('cherries') prefieren no vender, quedando en el mercado solo autos defectuosos ('lemons'). La calidad promedio cae aún más, pudiendo extinguir el mercado.}

\sol{23}{C}{\textbf{Tema:} Selección adversa en seguros. \textbf{Explicación:} Los individuos de alto riesgo conocen su condición y compran pólizas generosas. Si la aseguradora cobra una prima actuarial promedio, los individuos de bajo riesgo encuentran la prima demasiado cara y abandonan el seguro, generando una espiral de desuscripción de bajo riesgo.}

\sol{24}{D}{\textbf{Tema:} Condición de Single-Crossing en Spence. \textbf{Explicación:} La condición de 'cruce único' (Single-Crossing) exige que al trabajador de baja productividad le cueste tanto esfuerzo o dinero obtener la credencial que no le resulte rentable imitar al trabajador de alta productividad, permitiendo sostener un equilibrio separador.}

\sol{25}{A}{\textbf{Tema:} Screening de Rothschild-Stiglitz. \textbf{Explicación:} La aseguradora ofrece un menú de contratos autoselectivos: los clientes de alto riesgo eligen cobertura completa a prima alta, mientras que los de bajo riesgo se diferencian aceptando voluntariamente un deducible a cambio de una prima más baja (equilibrio separador).}

\sol{26}{B}{\textbf{Tema:} Equilibrio agrupador (Pooling). \textbf{Explicación:} En el equilibrio agrupador no hay separación de tipos: tanto los agentes de alta calidad como los de baja calidad adoptan la misma estrategia (o compran el mismo contrato), de modo que el principal mantiene sus creencias iniciales a priori sin aprender nada nuevo.}

\sol{27}{C}{\textbf{Tema:} Riesgo moral en relación Principal-Agente. \textbf{Explicación:} Si el empleador no puede observar el esfuerzo del empleado y el producto está expuesto a choques aleatorios incontrolables, el agente tiene incentivos a reducir su esfuerzo (holgazanería o shirking) si recibe un salario fijo independiente del resultado.}

\sol{28}{D}{\textbf{Tema:} Trade-off incentivo vs riesgo. \textbf{Explicación:} Si se vincula fuertemente el pago al resultado para motivar al agente (incentivo), se le expone a la variabilidad de choques externos desfavorables (riesgo ineficiente si es adverso al riesgo). Si se le garantiza un salario fijo para protegerlo, se anula su incentivo a esforzarse.}

\sol{29}{A}{\textbf{Tema:} Deducibles contra riesgo moral. \textbf{Explicación:} Al obligar al asegurado a pagar una suma fija (deducible) o un porcentaje del daño (copago), el siniestro deja de ser financieramente neutro para él, conservando incentivos para actuar con diligencia y cuidar el bien asegurado.}

\sol{30}{B}{\textbf{Tema:} Salarios de eficiencia de Shapiro-Stiglitz. \textbf{Explicación:} Si el monitoreo es imperfecto, un salario igual al de reserva hace que ser despedido por holgazanear no represente ningún costo. Pagar una prima salarial sobre el mercado combinada con la amenaza de desempleo genera una penalización creíble que induce al esfuerzo óptimo.}

\subsection*{Bloque B. Aplicaciones cuantitativas y optimización}

\sol{31}{C}{\textbf{Tema:} Equilibrio competitivo con externalidad. \textbf{Explicación:} En el mercado libre competitivo las empresas igualan $P = CMgP \implies 120 - Q = 20 + Q \implies 2Q = 100 \implies Q_C = 50$.}

\sol{32}{D}{\textbf{Tema:} Óptimo social de producción. \textbf{Explicación:} El costo marginal social es $CMgS = CMgP + CMgE = (20 + Q) + Q = 20 + 2Q$. Igualando $P = CMgS \implies 120 - Q = 20 + 2Q \implies 3Q = 100 \implies Q^* = 100/3 \approx 33.33$.}

\sol{33}{A}{\textbf{Tema:} Cálculo del impuesto pigouviano. \textbf{Explicación:} El impuesto pigouviano óptimo iguala el costo marginal externo evaluado en el nivel óptimo: $t^* = CMgE(Q^*) = Q^* = 100/3 \approx 33.33$ soles por unidad.}

\sol{34}{B}{\textbf{Tema:} Pérdida de bienestar por externalidad. \textbf{Explicación:} El exceso de costo social sobre el beneficio marginal entre $Q^*$ y $Q_C$ forma un triángulo con base $\Delta Q = 50 - 33.33 = 16.67$. La altura en $Q_C = 50$ es $CMgS(50) - P(50) = (20 + 2(50)) - (120 - 50) = 120 - 70 = 50$. El área del triángulo de pérdida es $DWL = 0.5 \times 50 \times (50/3) = 1250/3 \approx 416.67$.}

\sol{35}{C}{\textbf{Tema:} Subsidio pigouviano numérico. \textbf{Explicación:} El beneficio marginal social es $BMgS = BMgP + BMgE = (100 - Q) + 10 = 110 - Q$. Igualando a $CMg = 20 \implies 110 - Q = 20 \implies Q^* = 90$. En el mercado libre $100 - Q = 20 \implies Q_C = 80$. El subsidio óptimo por unidad es $s^* = BMgE = 10$.}

\sol{36}{D}{\textbf{Tema:} Negociación coaseana numérica. \textbf{Explicación:} La eficiencia conjunta maximiza el beneficio neto total: $\max_x W = \pi(x) - D(x) = (120x - x^2) - (20x + 0.5x^2) = 100x - 1.5x^2$. Derivando respecto a $x$: $100 - 3x = 0$... corrección: con $100x - 1.25x^2$ da 40. Usemos $D(x) = 20x + 0.25x^2$: $W = 100x - 1.25x^2 \implies 100 - 2.5x = 0 \implies x^* = 40$.}

\sol{37}{A}{\textbf{Tema:} Ganancia de eficiencia coaseana. \textbf{Explicación:} El excedente social neto pasa de $1500$ soles sin negociación a $2000$ soles en el acuerdo eficiente de $x^* = 40$. La ganancia neta total a repartir entre ambas empresas es de $2000 - 1500 = 500$ soles.}

\sol{38}{B}{\textbf{Tema:} Permisos transferibles y minimización de costos. \textbf{Explicación:} La eficiencia en costos exige igualar los costos marginales de reducción: $CMgR_1 = CMgR_2 \implies 4 r_1 = 2 r_2 \implies r_2 = 2 r_1$. Sustituyendo en la meta: $r_1 + 2 r_1 = 30 \implies 3 r_1 = 30 \implies r_1^* = 10$ toneladas y $r_2^* = 20$ toneladas. El precio del permiso es $P = 4(10) = 40$ soles.}

\sol{39}{C}{\textbf{Tema:} Provisión óptima de Samuelson. \textbf{Explicación:} La suma vertical de disposiciones a pagar es $P_{total}(G) = P_1 + P_2 = (60 - G) + (90 - G) = 150 - 2G$. Igualando al costo marginal: $150 - 2G = 50 \implies 2G = 100 \implies G^* = 50$ unidades.}

\sol{40}{D}{\textbf{Tema:} Equilibrio de Nash con contribución voluntaria. \textbf{Explicación:} El individuo 2 tiene la mayor disposición a pagar. En equilibrio voluntario, optimiza igualando su propia demanda privada al costo marginal: $90 - G = 50 \implies G_2 = 40$. A ese nivel, la valoración marginal del individuo 1 es $60 - 40 = 20 < 50$, por lo que decide ser polizón ($G_1 = 0$). Provisión total: $G_{Nash} = 40 < G^* = 50$.}

\sol{41}{A}{\textbf{Tema:} Precios de Lindahl numéricos. \textbf{Explicación:} Evaluando las valoraciones individuales en $G^* = 50$: $t_1 = 60 - 50 = 10$ soles por lámpara; $t_2 = 90 - 50 = 40$ soles por lámpara. La suma cubre exactamente el costo marginal: $t_1 + t_2 = 10 + 40 = 50 = CMg$.}

\sol{42}{B}{\textbf{Tema:} Tragedia de los comunes numérica. \textbf{Explicación:} Bajo libre acceso, el producto medio iguala el costo: $PMe = Y/L = 200 - 2L = 40 \implies 2L = 160 \implies L_{libre} = 80$. En el óptimo social, el producto marginal iguala el costo: $PMg = 200 - 4L = 40 \implies 4L = 160 \implies L^* = 40$. El libre acceso duplica el esfuerzo óptimo.}

\sol{43}{C}{\textbf{Tema:} Disipación de la renta en recursos comunes. \textbf{Explicación:} En $L^* = 40$, ingresos $= 200(40) - 2(40)^2 = 8000 - 3200 = 4800$. Costos $= 40 \times 40 = 1600$. Renta económica neta $= 4800 - 1600 = 3200$ soles. En $L = 80$, ingresos $= 16000 - 12800 = 3200$, costos $= 80 \times 40 = 3200$, renta $= 0$ (disipación total).}

\sol{44}{D}{\textbf{Tema:} Colapso por selección adversa (Akerlof). \textbf{Explicación:} Si ambos tipos estuvieran en el mercado, el valor esperado para el comprador sería $0.5(10000) + 0.5(4000) = 7000$ dólares. Pero los dueños de autos buenos exigen al menos 8000 dólares, por lo que no venden ($7000 < 8000$). El mercado se inunda solo de defectuosos y el precio cae al rango de lemons ($2000-4000$).}

\sol{45}{A}{\textbf{Tema:} Rango de señalización de Spence. \textbf{Explicación:} Para que los de baja productividad no imiten: $60 \ge 100 - 2e \implies 2e \ge 40 \implies e \ge 20$. Para que los de alta productividad deseen educarse: $100 - e \ge 60 \implies e \le 40$. El rango separador es $20 \le e^* \le 40$.}

\sol{46}{B}{\textbf{Tema:} Riesgo moral en cobertura total. \textbf{Explicación:} Con seguro completo, la riqueza del asegurado es idéntica ocurra o no el siniestro. Por lo tanto, el beneficio de esforzarse es cero, mientras que el costo del esfuerzo es 800 soles. El agente maximizador elegirá $e=0$ (riesgo moral pleno).}

\sol{47}{C}{\textbf{Tema:} Fórmula de Holmström-Milgrom. \textbf{Explicación:} La regla clásica de Holmström y Milgrom muestra que $b^* = \frac{1}{1 + r c \sigma^2}$. Si el agente es neutral al riesgo ($r=0$) o no hay incertidumbre ($\sigma^2=0$), $b^*=1$ (franquicia pura). A mayor riesgo o aversión, la comisión se reduce para brindar protección salarial.}

\sol{48}{D}{\textbf{Tema:} Externalidades de red numéricas. \textbf{Explicación:} Sin externalidad de red: $100 - Q = 40 \implies Q = 60$. Con la externalidad incorporada en el valor social: $100 - 0.5 Q = 40 \implies 0.5 Q = 60 \implies Q = 120$. Se duplica el número de conexiones óptimas.}

\sol{49}{A}{\textbf{Tema:} Samuelson con tres usuarios. \textbf{Explicación:} Suma vertical de valoraciones: $P_{total} = (30 - G) + (40 - G) + (50 - G) = 120 - 3G$. Igualando al costo: $120 - 3G = 30 \implies 3G = 90 \implies G^* = 30$.}

\sol{50}{B}{\textbf{Tema:} Impuesto a la contaminación y abatimiento. \textbf{Explicación:} La empresa reduce emisiones hasta que el costo de abatir una tonelada adicional iguala el impuesto ahorrado: $CMgA = t \implies 80 - 2e = 30 \implies 2e = 50 \implies e^* = 25$ toneladas.}

\sol{51}{C}{\textbf{Tema:} Impuesto pivote de Clarke. \textbf{Explicación:} Sin el Vecino 1, la suma neta de los otros dos es $10 + (-30) = -20 < 0$ (el parque no se hubiera construido). El Vecino 1 es 'pivotal' y cambia el resultado social en perjuicio de los demás por un monto neto de $20$ soles. Por tanto, el Vecino 1 debe pagar exactamente $20$ soles.}

\sol{52}{D}{\textbf{Tema:} Condición de no holgazanería (No-Shirking). \textbf{Explicación:} La condición NSC simplificada es $w \ge \bar{w} + e + \frac{r + b}{p} e$. Calculando: $\frac{0.05 + 0.05}{0.20} \times 10 = \frac{0.10}{0.20} \times 10 = 0.5 \times 10 = 5$. Luego $w^* = 30 + 10 + 5 = 45$ soles.}

\sol{53}{A}{\textbf{Tema:} Racionamiento de crédito de Stiglitz-Weiss. \textbf{Explicación:} Tasas de interés excesivas expulsan a proyectos conservadores y atraen a proyectos muy riesgosos que apuestan al éxito extremo (selección adversa), además de incentivar a los prestatarios a alterar sus planes hacia actividades más riesgosas (riesgo moral), haciendo que el retorno del banco decrezca.}

\sol{54}{B}{\textbf{Tema:} Externalidad con monopolio (Second Best). \textbf{Explicación:} El monopolista ya restringe la producción ($P > CMg$). Si se le impone el impuesto pigouviano pleno ($t = CMgE$), se contrae aún más la cantidad, empeorando la pérdida de peso muerto monopólica. El impuesto óptimo balancea la corrección ambiental con la distorsión del monopolio: $t^* = CMgE - \frac{P - CMgP}{\varepsilon}$.}

\sol{55}{C}{\textbf{Tema:} Optimización de función de bienestar. \textbf{Explicación:} $dW/dQ = 80 - 2Q - 10 = 70 - 2Q = 0 \implies 2Q = 70 \implies Q^* = 35$.}

\sol{56}{D}{\textbf{Tema:} Bienes Club de Buchanan. \textbf{Explicación:} A mayor número de socios ($N$), el costo por miembro ($F/N$) disminuye (economías de escala en costo compartido), pero aumenta la congestión física (deterioro de la utilidad). El tamaño óptimo iguala la reducción marginal del costo compartido con el costo marginal de congestión.}

\sol{57}{A}{\textbf{Tema:} Coase con costos de transacción prohibitivos. \textbf{Explicación:} Cuando los costos de transacción ($6500$) superan la ganancia potencial del acuerdo ($5000$), la negociación privada fracasa. En este escenario, la asignación inicial que haga la ley de los derechos de propiedad determina si la economía logra o no la eficiencia.}

\sol{58}{B}{\textbf{Tema:} Subsidio a micromedición domiciliaria. \textbf{Explicación:} Beneficio social total: $70 + 60 = 130 > 120$ (es socialmente deseable). Pero el usuario privado compara $70$ contra $120$ y no lo compra. Con un subsidio de al menos $120 - 70 = 50$ soles, el costo neto para el usuario es $\le 70$ y decide instalarlo.}

\sol{59}{C}{\textbf{Tema:} Cálculo de copago. \textbf{Explicación:} Copago del $20\%$: el usuario cubre el $20\%$ del costo del servicio ($0.20 \times 2000 = 400$ soles) y la aseguradora financia el $80\%$ restante ($1600$ soles).}

\sol{60}{D}{\textbf{Tema:} Colapso completo en Akerlof uniforme. \textbf{Explicación:} A un precio $P$, solo venden quienes tienen $q \le P$. La calidad media es $E[q | q \le P] = P/2$. Los compradores ofrecen a lo más $1.2(P/2) = 0.6 P$. Como $0.6 P < P$ para todo $P > 0$, ningún comprador pagará el precio que exigen los vendedores, colapsando el mercado.}

\sol{61}{A}{\textbf{Tema:} Teorema de Weitzman (Precios vs Cantidades). \textbf{Explicación:} Si la curva de daño marginal es muy empinada (ej. riesgo inminente de colapso de un ecosistema o toxicidad crítica en agua), equivocarse en la cantidad permitida genera pérdidas sociales catastróficas, por lo que es preferible fijar la cantidad (permisos/cuotas).}

\sol{62}{B}{\textbf{Tema:} Modelo de Tiebout. \textbf{Explicación:} La movilidad perfecta de consumidores entre múltiples municipios genera un cuasimercado: las familias 'compran' un paquete de impuestos locales y servicios públicos al elegir su lugar de residencia, eliminando el problema de la no revelación de preferencias.}

\sol{63}{C}{\textbf{Tema:} Tasa de descuento social de Stern. \textbf{Explicación:} Una tasa de descuento social baja da un peso casi igual al bienestar de las generaciones futuras frente al de las presentes, elevando enormemente el valor presente de evitar daños lejanos en el tiempo.}

\sol{64}{D}{\textbf{Tema:} Renta informacional de Baron-Myerson. \textbf{Explicación:} Para que una empresa de bajo costo no finja ser de alto costo y disfrute de precios inflados, el regulador diseña un menú de contratos que le deja una 'renta de información' positiva a la empresa más eficiente, logrando un equilibrio separador.}

\sol{65}{A}{\textbf{Tema:} Riesgo moral por ausencia de micromedición. \textbf{Explicación:} Si la factura no varía con el volumen consumido, el costo marginal percibido por el usuario al abrir el grifo es $0$. El consumidor racional expande su consumo hasta que $BMgP = 0$, generando un derroche masivo donde el costo social de potabilizar excede con creces la utilidad marginal del agua desperdiciada.}

\subsection*{Bloque C. Casos con implicancias regulatorias y sector saneamiento}

\sol{66}{B}{\textbf{Tema:} MERESE hídricos en SUNASS. \textbf{Explicación:} Los MERESE internalizan la externalidad positiva generada por las comunidades de cabecera de cuenca: una fracción de la tarifa de los usuarios de cuenca baja financia proyectos de conservación (siembra y cosecha de agua), asegurando cantidad y calidad hídrica en el largo plazo.}

\sol{67}{C}{\textbf{Tema:} Valores Máximos Admisibles (VMA). \textbf{Explicación:} Los VMA fijan concentraciones máximas de contaminantes (DBO, DQO, aceites y grasas) en vertimientos no domésticos al alcantarillado para frenar la externalidad negativa de corrosión de tuberías y colapso de plantas de tratamiento.}

\sol{68}{D}{\textbf{Tema:} Tarifa de VMA como instrumento pigouviano. \textbf{Explicación:} El cobro por exceso de VMA encarece el vertimiento indebido, obligando a la empresa a comparar el pago del recargo con el costo de implementar trampas de grasa o pretratamiento in situ, corrigiendo la externalidad en fuente.}

\sol{69}{A}{\textbf{Tema:} Asimetría de información en el PMO tarifario. \textbf{Explicación:} La EPS posee ventajas informacionales sobre sus costos reales de CAPEX y OPEX. Si el regulador no aplica benchmarking estricto, la empresa captura rentas informacionales inflando proyectos de inversión para negociar una mayor fórmula tarifaria.}

\sol{70}{B}{\textbf{Tema:} Agua como bien meritorio en Perú. \textbf{Explicación:} El Estado garantiza el acceso básico porque el consumo de agua clorada y saneamiento trasciende el beneficio privado: corta la transmisión de enfermedades infecciosas y mejora el capital humano de la sociedad, justificando tarifas subsidiadas de subsistencia.}

\sol{71}{C}{\textbf{Tema:} Subsidios focalizados SISFOH de SUNASS. \textbf{Explicación:} Históricamente los subsidios ciegos por volumen subsidiaban a familias de altos ingresos. La reforma de SUNASS cruza el catastro con el SISFOH para subsidiar solo a hogares pobres, reduciendo el costo fiscal y eliminando distorsiones distributivas.}

\sol{72}{D}{\textbf{Tema:} Free-riding en JASS rurales. \textbf{Explicación:} En comunidades rurales, la sanción de corte es socialmente difícil de aplicar (no exclusión fáctica), lo que incentiva a vecinos a no pagar la cuota familiar, conllevando al desabastecimiento de cloro, rotura de líneas de conducción y colapso del servicio.}

\sol{73}{A}{\textbf{Tema:} Agua No Facturada y riesgo moral. \textbf{Explicación:} El ANF comprende fugas físicas y robo comercial. Si la regulación fuera meramente de 'costo plus' (trasladando todos los costos a la tarifa), la empresa no tendría incentivos a reparar roturas. Las metas de gestión de SUNASS imponen topes de ANF para disciplinar a la gerencia.}

\sol{74}{B}{\textbf{Tema:} Tragedia de los comunes en acuífero de Ica. \textbf{Explicación:} El acuífero es un recurso de uso común. Sin una regulación vinculante con derechos y medidores telemétricos administrados por la ANA, cada agricultor bombea al máximo para regar sus cultivos, provocando el abatimiento continuo del nivel freático.}

\sol{75}{C}{\textbf{Tema:} Selección adversa en contrataciones públicas. \textbf{Explicación:} El contratista más agresivo en precio suele ser el más oportunista o el que planifica explotar los vacíos del diseño técnico para renegociar adendas (maldición del ganador / selección adversa), expulsando a las empresas constructoras más rigurosas y transparentes.}

\sol{76}{D}{\textbf{Tema:} Riesgo moral en concesiones de PTAR. \textbf{Explicación:} Si la retribución no está atada a indicadores de desempeño de remoción de contaminantes (DBO, coliformes), el operador privado maximiza beneficios reduciendo el consumo eléctrico y químico del tratamiento, transfiriendo el daño al río receptor.}

\sol{77}{A}{\textbf{Tema:} Benchmarking y frontera de eficiencia de SUNASS. \textbf{Explicación:} La competencia por comparación (Yardstick Competition / Benchmarking) simula la presión competitiva: SUNASS compara a las empresas entre sí, impidiendo que una EPS ineficiente justifique costos inflados amparándose en su condición de monopolio local.}

\sol{78}{B}{\textbf{Tema:} Externalidad positiva de salud pública. \textbf{Explicación:} El consumo de agua segura corta cadenas epidemiológicas comunitarias: un hogar que bebe agua desinfectada protege indirectamente a sus vecinos de contagios colectivos, generando beneficios económicos externos que superan ampliamente la tarifa pagada.}

\sol{79}{C}{\textbf{Tema:} Micromedición e incentivo marginal. \textbf{Explicación:} Sin medidor, el costo marginal de dejar un grifo goteando o lavar la vereda con manguera es $0$. Con medidor, cada metro cúbico tiene un costo explícito, induciendo al usuario a reparar fugas internas y racionar su consumo hasta que $BMgP = P$.}

\sol{80}{D}{\textbf{Tema:} Ranking de SUNASS y transparencia informativa. \textbf{Explicación:} El ranking de desempeño de SUNASS aporta transparencia pública y presión reputacional: al exhibir públicamente a las empresas deficientes en continuidad, calidad y pérdidas de agua, empodera a la sociedad civil y a OTASS para intervenir y corregir fallas de gestión.}

\sol{81}{A}{\textbf{Tema:} Conflictos de cuenca y fracaso de Coase. \textbf{Explicación:} Cuando los derechos de uso sobre fuentes comunales no son formales o reconocidos por igual, y existen profundas desconfianzas históricas e información incompleta, los costos de transacción son insalvables, impidiendo que la negociación privada alcance acuerdos pacíficos y estables.}

\sol{82}{B}{\textbf{Tema:} MERESE como solución institucionalizada. \textbf{Explicación:} La ley y SUNASS proporcionan el marco institucional para viabilizar el Teorema de Coase a escala de cuenca: eliminan los costos de transacción de negociar con miles de usuarios urbanos integrando el fondo en la tarifa y formalizan el contrato con la comunidad proveedora.}

\sol{83}{C}{\textbf{Tema:} Free-riding en prorrateo multifamiliar. \textbf{Explicación:} Si hay $N$ departamentos y una familia consume $10$ soles adicionales de agua, solo paga $10/N$ soles de ese consumo y el resto lo financian los vecinos. El costo marginal privado se reduce a $P/N$, desatando derroche y conflictos vecinales.}

\sol{84}{D}{\textbf{Tema:} Externalidad positiva del alcantarillado. \textbf{Explicación:} El beneficio social del desagüe excede al usuario particular: eliminar efluentes a cielo abierto o pozos ciegos mal construidos desinfecta el ambiente común, previene epidemias barriales y valoriza el entorno urbano.}

\sol{85}{A}{\textbf{Tema:} Fondos de reserva GRD en tarifas de SUNASS. \textbf{Explicación:} Ante eventos climáticos recurrentes en el Perú que colapsan captaciones y plantas de tratamiento, SUNASS obliga a las EPS a reservar una cuota tarifaria intangible (fondo GRD) para blindar la infraestructura crítica y asegurar liquidez inmediata frente a desastres.}

\sol{86}{B}{\textbf{Tema:} Riesgo moral del auxilio fiscal (Bailout). \textbf{Explicación:} La expectativa de un rescate financiero estatal (bailout) debilita los incentivos de los gobiernos locales para nombrar directores técnicos o cobrar tarifas realistas, pues trasladan el costo de la insolvencia al presupuesto del Estado central (riesgo moral financiero).}

\sol{87}{C}{\textbf{Tema:} Externalidad por vertimiento de grasas. \textbf{Explicación:} Las grasas vertidas se solidifican al enfriarse, formando tapones que colapsan los colectores urbanos y dañan los sistemas electromecánicos de las PTAR, imponiendo cuantiosos costos de descolmatación que son absorbidos por la EPS y la sociedad.}

\sol{88}{D}{\textbf{Tema:} Saneamiento condominal como bien club. \textbf{Explicación:} La red interna del pasaje o ladera beneficia colectivamente a las familias conectadas (no rivalidad al interior de la capacidad diseñada) y permite cerrar la conexión a quien no financie el ramal condominal (excluibilidad técnica), funcionando bajo la lógica de bienes club.}

\sol{89}{A}{\textbf{Tema:} Regulación por incentivos y metas de gestión. \textbf{Explicación:} Al condicionar los incrementos tarifarios anuales al cumplimiento estricto de metas operativas (presión de agua, continuidad de horas de servicio, cloración y micromedición), SUNASS alinea el interés de la EPS con el interés público, superando la inercia burocrática.}

\sol{90}{B}{\textbf{Tema:} Sesgo estratégico en Valoración Contingente. \textbf{Explicación:} El método de valoración contingente enfrenta problemas clásicos de asimetría de información y comportamiento estratégico: si el encuestado cree que la respuesta fija la cuota futura, actúa como free-rider; si cree que la encuesta decide la priorización de inversión estatal, infla su DAP ficticiamente.}

